% !TeX program = xelatex
\documentclass[11pt,a4paper]{article}

\usepackage{fontspec}
\setmainfont{Libertinus Serif}
\usepackage{polyglossia}
\setdefaultlanguage{serbian}
\setotherlanguage{english}
\usepackage{geometry}
\usepackage{graphicx}
\usepackage{xcolor}
\usepackage{enumitem}
\usepackage{hyperref}
\usepackage{bookmark}
\usepackage{microtype}
\usepackage{parskip}
\usepackage{fvextra}
\usepackage{float}
\usepackage{caption}

\geometry{margin=2.2cm}
\setlength{\emergencystretch}{3em}
\setlist{leftmargin=*}
\urlstyle{same}
\hypersetup{
  colorlinks=true,
  linkcolor=blue!50!black,
  urlcolor=blue!50!black,
  pdfcreator={LaTeX via pandoc}
}
\providecommand{\tightlist}{%%
  \setlength{\itemsep}{0pt}\setlength{\parskip}{0pt}}
\DefineVerbatimEnvironment{CodeBlock}{Verbatim}{breaklines=true,breakanywhere=true,fontsize=\small}

\begin{document}
\section{C\# - основе}

Уз овај документ ћеш проћи кроз основе \emph{C\#} програмског језика као
и његову употребу унутар \emph{Visual Studio Code}-а. У учионицама у
којима се одржавају вјежбе имате двије опције за израду данашњих
задатака: једна је да користите виртуелну машину и радите унутар
\emph{Visual Studio}-а окружења, а друга је да користите \emph{Visual
Studio Code} који је доступан изван виртуелне машине. Топла је препорука
да користите онај доступан унутар виртуелне машине, као и да на својим
персоналним рачунарима не користите \emph{VS Code}, већ неко од
интегрисаних окружења за развој (IDE-a):
\href{https://visualstudio.microsoft.com/}{\emph{Visual Studio}} и
\href{https://www.jetbrains.com/rider/}{\emph{JetBrains Rider}} су
најпознатији (за оба алата можеш добити професионалне лиценце на основу
факултетског имејла).

Уколико ћеш радити у \emph{Code}-у ће ти требати одговарајући
\emph{plugin} за \emph{C\#}. Препорука је да у фолдеру \emph{Documents}
направиш фолдер са својим именом и индексом и све задатке смјестиш у
њега. Исте налоге користе сви студенти, па је ово препоручен начин за
рад да не изгубиш задатке од претходног пута. За сваки од задатака у
наставку направи нови фолдер унутар твог фолдера и у том фолдеру направи
нову апликацију у којој ћеш ријешити задатак.

Уколико ћеш радити у \emph{Visual Studio}-у, немаш никаквих предуслова.

За сваки задатак направи нови пројекат/нову апликацију. За оне који желе
да науче више:
\href{https://learn.microsoft.com/en-us/previous-versions/visualstudio/visual-studio-2015/ide/adding-and-removing-project-items?view=vs-2015}{погледајте}
како да унутар једне апликације имате више пројеката.

\section{Развој конзолних
апликација}

За израду задатака који слиједе може ти значити
\href{https://www.codecademy.com/learn/learn-c-sharp/modules/csharp-hello-world/cheatsheet}{овај}
и
\href{https://begincodingnow.com/wp-content/uploads/2019/02/CSharpCheatSheetV1.pdf}{овај}
(мало старија верзија) документ, као и званичну
\href{https://learn.microsoft.com/en-us/dotnet/csharp/language-reference/}{документацију}.

\begin{enumerate}
\def\labelenumi{\arabic{enumi}.}
\item
  Пратећи упутство са
  \href{https://learn.microsoft.com/en-us/previous-versions/visualstudio/visual-studio-2015/ide/creating-solutions-and-projects?view=vs-2015\#create-a-project-from-an-installed-project-template}{VS}/
  \href{https://docs.microsoft.com/en-us/dotnet/core/tutorials/with-visual-studio-code?pivots=dotnet-6-0\#create-the-app}{VS
  Code}, направи своју прву конзолну апликацију (енгл. \emph{Console
  Application}). Уколико радите у \emph{Code}-у, обратите пажњу на тачку
  7 из упутства и користите ту синтаксу.
\item
  Направи још једну конзолну апликацију која учитава природан број
  \emph{n} и као резултат испише \emph{n}-ти члан Фибоначијевог низа.
\item
  Прочитај документ о
  \href{https://docs.microsoft.com/en-us/dotnet/csharp/fundamentals/coding-style/coding-conventions}{препорученим
  конвенција}ма. Анализирај колико си испоштовао конвенције у претходна
  два задатка. Рефакториши претходне задатке тако да испоштујеш
  конвенције и труди се да их се придржавате убудуће.
\item
  Направи конзолну апликацију која учитава стринг и провјерава да ли је
  стринг палиндром.
\item
  Направи конзолну апликацију која учитава стринг произвољне дужине и
  серијализује га у текстуални фајл.
\end{enumerate}

\subsection{Класе}

\begin{enumerate}
\def\labelenumi{\arabic{enumi}.}
\setcounter{enumi}{5}
\item
  Направи нову конзолну апликацију у којој ћеш вјежбати рад са класама.
  Направи нову класу аутомобил која има поља: модел, марка, боја, година
  производње, врсту горива и мотор. Врста горива је енумерација (дизел и
  бензин). Мотор је такође класа са идућим пољима: назив, кубикажа,
  снага.
\item
  Направи пар објеката класа мотор и аутомобил. За обје класе направи
  \emph{ToString} методу и испиши објекте на конзолу.
\end{enumerate}

\subsection{Инсталација и коришћење
пакета}

\begin{enumerate}
\def\labelenumi{\arabic{enumi}.}
\setcounter{enumi}{7}
\item
  Пратећи упутство са
  \href{https://learn.microsoft.com/en-us/nuget/consume-packages/install-use-packages-visual-studio}{VS}/\href{https://docs.microsoft.com/en-us/nuget/quickstart/install-and-use-a-package-using-the-dotnet-cli}{VS
  Code}, инсталирај \emph{Newtonsoft.Json} пакет и тестирајте га на
  примјеру.
\item
  Вратите се на пројекат који си користио за вјежбање рада са класама. У
  новој функцији учитај \emph{cars.json} фајл (налази се на канвасу у
  фолдеру вјежбе/термин 1) користећи \emph{Newtonsoft.Json} пакет
  (користи ово подешавање приликом учитавања new JsonSerializerSettings
  \{ PreserveReferencesHandling = PreserveReferencesHandling.Objects \},
  а у документацији пакета пронађи како да га искористиш). На конзолу
  испиши:

  \begin{enumerate}
  \def\labelenumii{\alph{enumii}.}
  \item
    Старост сваког аутомобила
  \item
    Називе аутомобила који су највише 5 година стари
  \item
    Просјечну старост
  \end{enumerate}
\item
  Из листе аутомобила издвоји све јединствене моторе и серијализуј их у
  фајл под називом \emph{engines.json}.
\item
  У новој конзолној апликацији направи класу производ и учитај
  \emph{items.json}. Испиши све производе коjима рок трајања није
  истекао, па испиши оне којима истиче у року од 7 дана.
\item
  Производе из претходног задатка пребаци у ријечник тако да су кључеви
  шифре производа, а вриједности су сами производи. Користећи ријечник
  израчунај колико ће корисника коштати куповина два производа са шифром
  1, једног производа са шифром 3 и четири производа са шифром 4.
\end{enumerate}

\end{document}
